\documentclass[11pt]{article}
\usepackage[margin=1in]{geometry}
\usepackage{amsmath,amssymb,bm,braket}
\usepackage{hyperref}
\begin{document}
Let

\[
\omega=e^{2\pi i/N},
\]

and use the standard parafermion algebra

\[
\alpha_i^N=1,
\qquad
\alpha_i\alpha_j=\omega\,\alpha_j\alpha_i
\quad (i<j).
\]

This convention agrees with the usual parity-operator formulation of $\mathbb Z_N$ parafermion braiding; see \href{\detokenize{https://arxiv.org/html/1511.02704v1}}{Hutter and Loss}, particularly their Eqs. (1), (6), and (23). Adiabatic exchange acts as a non-Abelian geometric phase in the degenerate ground-state manifold, as discussed for parafermion tunneling protocols in \href{\detokenize{https://arxiv.org/abs/2403.09602}}{Hong et al.}.

\paragraph*{Step-by-Step Derivation}

\paragraph*{1. Fusion-parity operators}

For $i<j$, define

\[
Q_{ij}
=
\omega^{(N+1)/2}\alpha_i\alpha_j^\dagger,
\qquad
Q_{ij}^N=1.
\]

A fusion state with channel $p\in\mathbb Z_N$ satisfies

\[
Q_{ij}|p\rangle=\omega^p|p\rangle.
\]

Using $\alpha_j\alpha_i^\dagger=\omega\alpha_i^\dagger\alpha_j$,

\[
Q_{ij}^\dagger
=
\omega^{-(N+1)/2}\alpha_j\alpha_i^\dagger
=
\omega^{(1-N)/2}\alpha_i^\dagger\alpha_j,
\]

so

\[
\alpha_i^\dagger\alpha_j
=
\omega^{(N-1)/2}Q_{ij}^\dagger.
\]

Consequently, in the $Q_{ij}=\omega^p$ sector,

\[
E_p
=
2t\cos\!\left[
\frac{\pi(N-1)-\phi_{ij}-2\pi p}{N}
\right].
\]

Write

\[
-\frac{\phi_{ij}}{2\pi}=k_{ij}+\delta,
\qquad 0<\delta<1.
\]

For $p=k_{ij}$,

\[
E_{k_{ij}}
=
2t\cos\!\left[
\frac{\pi(N-1+2\delta)}{N}
\right],
\]

which is the level closest to the minimum of the cosine. Hence the ground-state projector of $H_{ij}$ is

\[
P_{ij}^{(k_{ij})}
=
\frac1N\sum_{n=0}^{N-1}
\omega^{-k_{ij}n}Q_{ij}^{\,n}.
\]

Thus $k_{ij}$ is precisely the ground-state fusion channel stated in the problem.

\paragraph*{2. Convenient basis}

Choose basis states $|x,y\rangle$ satisfying

\[
Q_{12}|x,y\rangle=\omega^x|x,y\rangle,
\qquad
Q_{34}|x,y\rangle=\omega^y|x,y\rangle.
\]

The initial state is therefore

\[
|\psi^i(q)\rangle=|q,k_{34}\rangle.
\]

The relevant parity algebra is

\[
Q_{12}Q_{23}=\omega Q_{23}Q_{12},
\]

and one may choose the phases of $|x,y\rangle$ so that

\[
Q_{23}|x,y\rangle=|x+1,y-1\rangle.
\]

Moreover,

\[
Q_{13}
=
\omega^{(N-1)/2}Q_{12}Q_{23},
\]

and hence

\[
Q_{13}|x,y\rangle
=
\omega^{x+(N+1)/2}|x+1,y-1\rangle.
\]

Iterating this action gives

\[
Q_{13}^{m}|a,y\rangle
=
\omega^{ma+\frac{m(m+N)}2}
|a+m,y-m\rangle.
\]

All channel labels and additions are understood modulo $N$.

\paragraph*{3. Apply the tunneling sequence}

For compactness, set

\[
a=k_{12},\qquad
b=k_{13},\qquad
c=k_{23},\qquad
d=k_{34}.
\]

The normalized ground-space transport operator is proportional to

\[
U
=
N^2P_{34}^{(d)}
P_{13}^{(b)}
P_{12}^{(a)}
P_{23}^{(c)}
P_{34}^{(d)}.
\]

The factor $N^2$ removes the magnitude $N^{-2}$ from the four mutually unbiased ground-space projections.

First,

\[
P_{23}^{(c)}|q,d\rangle
=
\frac1N
\sum_{n=0}^{N-1}
\omega^{-cn}|q+n,d-n\rangle.
\]

Applying $P_{12}^{(a)}$ selects

\[
n=a-q.
\]

Therefore,

\[
P_{12}^{(a)}P_{23}^{(c)}|q,d\rangle
=
\frac1N
\omega^{-c(a-q)}
|a,d-a+q\rangle.
\]

Next, in $P_{13}^{(b)}$, the final $P_{34}^{(d)}$ selects the power

\[
m=q-a,
\]

because $Q_{13}^m$ changes

\[
|a,d-a+q\rangle
\longrightarrow
|a+m,d-a+q-m\rangle
=
|q,d\rangle.
\]

The resulting unnormalized amplitude is

\[
\frac{1}{N^2}
\omega^{-c(a-q)}
\omega^{-bm}
\omega^{ma+\frac{m(m+N)}2}.
\]

Since $m=q-a$ and $-c(a-q)=cm$, its exponent is

\[
\begin{aligned}
E(q)
&=
m(c-b+a)+\frac{m(m+N)}2\\
&=
\frac{q-a}{2}
\left[q+a+N+2(c-b)\right].
\end{aligned}
\]

Thus the normalized braid eigenvalue is

\[
e^{i\theta(q)}
=
\omega^{E(q)}
=
\exp\!\left[
\frac{i\pi}{N}
(q-k_{12})
\left(
q+k_{12}+N+2k_{23}-2k_{13}
\right)
\right].
\]

Equivalently, completing the square,

\[
e^{i\theta(q)}
=
\exp\!\left\{
\frac{i\pi}{N}
\left[
\left(q+k_{23}-k_{13}+\frac N2\right)^2
-
\left(k_{12}+k_{23}-k_{13}+\frac N2\right)^2
\right]
\right\}.
\]

The result is independent of $k_{34}$: that channel fixes the initially paired modes, but cancels because $Q_{23}$ and $Q_{13}$ conserve the total fusion channel $x+y$. Similarly, $t$ only controls the excitation gap and not the universal braiding phase.

An unspecified interpolation time can add a common, $q$-independent dynamical phase. The expression above is the fusion-channel-dependent geometric phase, normalized so that $\theta(k_{12})=0$. Reversing the braid orientation complex-conjugates it.

\textbf{Final Answer:}

\[
\boxed{
|\psi^f(q)\rangle
=
\exp\!\left[
\frac{i\pi}{N}
(q-k_{12})
\left(
q+k_{12}+N+2k_{23}-2k_{13}
\right)
\right]
|\psi^i(q)\rangle
}
\]

or, equivalently,

\[
\boxed{
|\psi^f(q)\rangle
=
\exp\!\left\{
\frac{i\pi}{N}
\left[
\left(q+k_{23}-k_{13}+\frac N2\right)^2
-
\left(k_{12}+k_{23}-k_{13}+\frac N2\right)^2
\right]
\right\}
|\psi^i(q)\rangle .
}
\]
\end{document}
